For $p=q=0.6$, $u=0.3$, $v=-0.2$, the complete physical kernel
on $a+c=b+d$ has the following pointwise reference values:
\begin{center}
\begin{tabular}{rrrrr}\toprule
$M$ & 2 & 4 & 8 & 16\\\midrule
$U_{00}$ & 0.96379723 & 1.15801422 & 1.36981862 & 1.58731944\\
\bottomrule\end{tabular}
\end{center}
A separate $M=32$ reference gives $1.80680785$.
These are point values, not a mass-integrated Bose observable.
The initial study also retains offsets $r=0.1,0.4$ at $M=2,4,8,16$.
No full-kernel limiting profile is inferred from these samples.

For the complete $M=32$ kernel, merely increasing arithmetic precision
does not repair an inadequate asymptotic tail. At 110 decimal digits,
the merged evaluator with $(N,K)=(512,24)$ has relative error about
$2.42\,10^{24}$ against the separate LPRS reference. The successive
settings $(1024,32)$, $(2048,40)$, $(4096,48)$ and $(8192,56)$ give
approximately $5.78\,10^6$, $1.74\,10^{-13}$, $5.39\,10^{-19}$ and
$8.60\,10^{-25}$. All attempts are retained. These are sample-specific
empirical errors, not a uniform accuracy claim for the algorithm.

For the single controlled Omega factor, the rigorous direct-sum
tail at $M=32$, $r=0$, $N=128$ is about $2.62\,10^{-39}$ in
exact arithmetic. At much larger $M$ this tail becomes far smaller
than the available rounding precision; it must not be reported as
the total numerical error. Computations at 70 and 100 decimal digits
are stored separately for the same series cutoff.

Convergence in \eqref{scaling} is slow: at $M=1024$ the complex scaled
value for $\lambda=0$ is approximately $1.31857+0.08411i$,
whose absolute difference from one is $0.32949$.
For $\lambda=1$ it is approximately $1.27753+0.09388i$,
with absolute difference $0.33222$ from $\sin(1/2)/(1/2)$.
The existence of the asymptotic limit does not imply these masses
are already close to it.

